% GENERATED FILE. Edit canonical lessons, then run npm run generate:books.
% canonical-source-hash: fnv1a64:8cf558ef605a2f88
% canonical-lessons: JA-C467-nuu, JA-C467-amimono, JA-C467-someru, JA-C467-shishuu, JA-C467-kezuru, JA-R467-review-462-467

\chapter{Sewing and Knitting}
\label{ch:ja-sewing-and-knitting}
\hlchaptermodality{467}

\begin{chapteropening}
\textbf{By the end of this chapter:} \emph{I can talk about things people make.}

Complete the last lesson of chapter 467: I can talk about things people make.
\end{chapteropening}

\section[\texorpdfstring{\ja{ぬう} (nuu)}{ぬう (nuu)}]{\ja{ぬう} (nuu) — to sew}

\label{lesson:JA-C467-nuu}



\begin{quote}
Before the new one: say the Japanese for paper folding, then the Japanese for the arts.
\end{quote}

\subsection*{You'll want to know: \ja{ぬう}}
\textbf{\ja{ぬう}} — \emph{nuu} — \textquotedblleft{}to sew\textquotedblright{}.

\begin{quote}
\textbf{\ja{ははは} \ja{ゆかたを} \ja{ぬっています。}} — \emph{haha wa yukata o nutte imasu} — \textquotedblleft{}My mother is sewing a cotton kimono.\textquotedblright{}
\end{quote}

Polite \textbf{\ja{ぬいます}}, \emph{nuimasu}; a godan verb. Past (polite): \textbf{\ja{そぼが} \ja{きものを} \ja{ぬいました。}}, \emph{sobo ga kimono o nuimashita}.

\begin{cousinweb}
\textbf{\ja{ぬう}} (\emph{nuu}) is a native Japanese word (wago), inherited rather than borrowed from Chinese; its kanji spelling was attached to it later (JMdict).
\end{cousinweb}

\begin{grammarlens}[title={What you've built}]
One more word for describing what you lived and felt.
\end{grammarlens}

\subsection*{Your turn}
\begin{itemize}
\raggedright
  \item \emph{Say it:} \emph{nuu}
  \item \emph{Say it:} \emph{nuu}, once more
  \item \emph{Say it:} \emph{haha wa yukata o nutte imasu} — My mother is sewing a cotton kimono.
  \item \emph{Recall:} say the Japanese for paper folding, then the Japanese for the arts, then say \emph{nuu} again
\end{itemize}

\begin{tcolorbox}[breakable,colback=teal!4,colframe=teal!35!black,title={Before you move on}]
What does \ja{ぬう} mean? (\textquotedblleft{}To sew\textquotedblright{}.) Say it once more.
\end{tcolorbox}
\section[\texorpdfstring{\ja{あみもの} (amimono)}{あみもの (amimono)}]{\ja{あみもの} (amimono) — knitting}

\label{lesson:JA-C467-amimono}



\begin{quote}
Before the new one: say the Japanese for the arts, then the Japanese for to sew.
\end{quote}

\subsection*{You'll want to know: \ja{あみもの}}
\textbf{\ja{あみもの}} — \emph{amimono} — \textquotedblleft{}knitting\textquotedblright{}.

\begin{quote}
\textbf{\ja{ふゆは} \ja{あみものを} \ja{します。}} — \emph{fuyu wa amimono o shimasu} — \textquotedblleft{}I do knitting in winter.\textquotedblright{}
\end{quote}

\begin{cousinweb}
\textbf{\ja{あみもの}} (\emph{amimono}) is a native Japanese compound (wago), inherited rather than borrowed from Chinese: \emph{ami}, \textquotedblleft{}knitting (amu)\textquotedblright{} and \emph{mono}, \textquotedblleft{}thing\textquotedblright{} (JMdict, KANJIDIC2).
\end{cousinweb}

\begin{grammarlens}[title={What you've built}]
One more word for describing what you lived and felt.
\end{grammarlens}

\subsection*{Your turn}
\begin{itemize}
\raggedright
  \item \emph{Say it:} \emph{amimono}
  \item \emph{Say it:} \emph{amimono}, once more
  \item \emph{Say it:} \emph{fuyu wa amimono o shimasu} — I do knitting in winter.
  \item \emph{Recall:} say the Japanese for the arts, then the Japanese for to sew, then say \emph{amimono} again
\end{itemize}

\begin{tcolorbox}[breakable,colback=teal!4,colframe=teal!35!black,title={Before you move on}]
What does \ja{あみもの} mean? (\textquotedblleft{}Knitting\textquotedblright{}.) Say it once more.
\end{tcolorbox}
\section[\texorpdfstring{\ja{そめる} (someru)}{そめる (someru)}]{\ja{そめる} (someru) — to dye}

\label{lesson:JA-C467-someru}



\begin{quote}
Before the new one: say the Japanese for to sew, then the Japanese for knitting.
\end{quote}

\subsection*{You'll want to know: \ja{そめる}}
\textbf{\ja{そめる}} — \emph{someru} — \textquotedblleft{}to dye\textquotedblright{}.

\begin{quote}
\textbf{\ja{ぬのを} \ja{あおく} \ja{そめました。}} — \emph{nuno o aoku somemashita} — \textquotedblleft{}I dyed the cloth blue.\textquotedblright{}
\end{quote}

Polite \textbf{\ja{そめます}}, \emph{somemasu}; an ichidan verb. Past (polite): \textbf{\ja{しょくにんが} \ja{ぬのを} \ja{そめました。}}, \emph{shokunin ga nuno o somemashita}.

\begin{cousinweb}
\textbf{\ja{そめる}} (\emph{someru}) is a native Japanese word (wago), inherited rather than borrowed from Chinese; its kanji spelling was attached to it later (JMdict).
\end{cousinweb}

\begin{grammarlens}[title={What you've built}]
One more word for describing what you lived and felt.
\end{grammarlens}

\subsection*{Your turn}
\begin{itemize}
\raggedright
  \item \emph{Say it:} \emph{someru}
  \item \emph{Say it:} \emph{someru}, once more
  \item \emph{Say it:} \emph{nuno o aoku somemashita} — I dyed the cloth blue.
  \item \emph{Recall:} say the Japanese for to sew, then the Japanese for knitting, then say \emph{someru} again
\end{itemize}

\begin{tcolorbox}[breakable,colback=teal!4,colframe=teal!35!black,title={Before you move on}]
What does \ja{そめる} mean? (\textquotedblleft{}To dye\textquotedblright{}.) Say it once more.
\end{tcolorbox}
\section[\texorpdfstring{\ja{ししゅう} (shishū)}{ししゅう (shishū)}]{\ja{ししゅう} (shishū) — embroidery}

\label{lesson:JA-C467-shishuu}



\begin{quote}
Before the new one: say the Japanese for knitting, then the Japanese for to dye.
\end{quote}

\subsection*{You'll want to know: \ja{ししゅう}}
\textbf{\ja{ししゅう}} — \emph{shishū} — \textquotedblleft{}embroidery\textquotedblright{}.

\begin{quote}
\textbf{\ja{ははの} \ja{ししゅうは} \ja{きれいです。}} — \emph{haha no shishū wa kirei desu} — \textquotedblleft{}My mother's embroidery is beautiful.\textquotedblright{}
\end{quote}

\begin{cousinweb}
\textbf{\ja{ししゅう}} (\emph{shishū}) is Sino-Japanese (kango): two Chinese characters read with their on'yomi, readings borrowed from Chinese. They are \emph{shi}, \textquotedblleft{}prick\textquotedblright{} (Middle Chinese \emph{tshjeH}) and \emph{shū}, \textquotedblleft{}sew\textquotedblright{} (\emph{sjuwH}) (KANJIDIC2; Middle Chinese in Baxter's transcription).
\end{cousinweb}

\begin{grammarlens}[title={What you've built}]
One more word for describing what you lived and felt.
\end{grammarlens}

\subsection*{Your turn}
\begin{itemize}
\raggedright
  \item \emph{Say it:} \emph{shishū}
  \item \emph{Say it:} \emph{shishū}, once more
  \item \emph{Say it:} \emph{haha no shishū wa kirei desu} — My mother's embroidery is beautiful.
  \item \emph{Recall:} say the Japanese for knitting, then the Japanese for to dye, then say \emph{shishū} again
\end{itemize}

\begin{tcolorbox}[breakable,colback=teal!4,colframe=teal!35!black,title={Before you move on}]
What does \ja{ししゅう} mean? (\textquotedblleft{}Embroidery\textquotedblright{}.) Say it once more.
\end{tcolorbox}
\section[\texorpdfstring{\ja{けずる} (kezuru)}{けずる (kezuru)}]{\ja{けずる} (kezuru) — to whittle; to shave down}

\label{lesson:JA-C467-kezuru}



\begin{quote}
Before the new one: say the Japanese for to dye, then the Japanese for embroidery.
\end{quote}

\subsection*{You'll want to know: \ja{けずる}}
\textbf{\ja{けずる}} — \emph{kezuru} — \textquotedblleft{}to whittle; to shave down\textquotedblright{}.

\begin{quote}
\textbf{\ja{ちちは} \ja{いたを} \ja{けずりました。}} — \emph{chichi wa ita o kezurimashita} — \textquotedblleft{}My father planed the board.\textquotedblright{}
\end{quote}

Polite \textbf{\ja{けずります}}, \emph{kezurimasu}; a godan verb. Past (polite): \textbf{\ja{だいくが} \ja{いたを} \ja{けずりました。}}, \emph{daiku ga ita o kezurimashita}.

\begin{cousinweb}
\textbf{\ja{けずる}} (\emph{kezuru}) is a native Japanese word (wago), inherited rather than borrowed from Chinese; its kanji spelling was attached to it later (JMdict).
\end{cousinweb}

\begin{grammarlens}[title={What you've built}]
One more word for describing what you lived and felt.
\end{grammarlens}

\subsection*{Your turn}
\begin{itemize}
\raggedright
  \item \emph{Say it:} \emph{kezuru}
  \item \emph{Say it:} \emph{kezuru}, once more
  \item \emph{Say it:} \emph{chichi wa ita o kezurimashita} — My father planed the board.
  \item \emph{Recall:} say the Japanese for to dye, then the Japanese for embroidery, then say \emph{kezuru} again
\end{itemize}

\begin{tcolorbox}[breakable,colback=teal!4,colframe=teal!35!black,title={Before you move on}]
What does \ja{けずる} mean? (\textquotedblleft{}To whittle; to shave down\textquotedblright{}.) Say it once more.
\end{tcolorbox}
\section[(dialogue)]{Looking back: from Win, Lose or Tie to Sewing and Knitting}

\label{lesson:JA-R467-review-462-467}



\begin{quote}
Say the Japanese for a contest and a tie.
\end{quote}

\subsection*{Your turn}
Say these aloud:

\begin{itemize}
\raggedright
  \item a contest, a tie, competition, a ranking, a foul
  \item to toughen, stamina, to aim at, to bounce, to vie
  \item spectators, a referee, points scored, to attack, to snatch
  \item paints, a brush, to paint a picture, a painter, fine art
  \item calligraphy, flower arranging, the tea ceremony, paper folding, the arts
  \item to sew, knitting, to dye, embroidery, to whittle
\end{itemize}

\begin{tcolorbox}[breakable,colback=teal!4,colframe=teal!35!black,title={Before you move on}]
A contest? (\textbf{\ja{しょうぶ}}, \emph{shōbu}.) A tie? (\textbf{\ja{ひきわけ}}, \emph{hikiwake}.) Competition? (\textbf{\ja{きょうそう}}, \emph{kyōsō}.) A ranking? (\textbf{\ja{じゅんい}}, \emph{jun'i}.) A foul? (\textbf{\ja{はんそく}}, \emph{hansoku}.) To toughen? (\textbf{\ja{きたえる}}, \emph{kitaeru}.) Stamina? (\textbf{\ja{たいりょく}}, \emph{tairyoku}.) To aim at? (\textbf{\ja{ねらう}}, \emph{nerau}.) To bounce? (\textbf{\ja{はずむ}}, \emph{hazumu}.) To vie? (\textbf{\ja{きそう}}, \emph{kisou}.) Spectators? (\textbf{\ja{かんきゃく}}, \emph{kankyaku}.) A referee? (\textbf{\ja{しんぱん}}, \emph{shinpan}.) Points scored? (\textbf{\ja{とくてん}}, \emph{tokuten}.) To attack? (\textbf{\ja{せめる}}, \emph{semeru}.) To snatch? (\textbf{\ja{うばう}}, \emph{ubau}.) Paints? (\textbf{\ja{えのぐ}}, \emph{enogu}.) A brush? (\textbf{\ja{ふで}}, \emph{fude}.) To paint a picture? (\textbf{\ja{えがく}}, \emph{egaku}.) A painter? (\textbf{\ja{がか}}, \emph{gaka}.) Fine art? (\textbf{\ja{びじゅつ}}, \emph{bijutsu}.) Calligraphy? (\textbf{\ja{しょどう}}, \emph{shodō}.) Flower arranging? (\textbf{\ja{いけばな}}, \emph{ikebana}.) The tea ceremony? (\textbf{\ja{さどう}}, \emph{sadō}.) Paper folding? (\textbf{\ja{おりがみ}}, \emph{origami}.) The arts? (\textbf{\ja{げいじゅつ}}, \emph{geijutsu}.) To sew? (\textbf{\ja{ぬう}}, \emph{nuu}.) Knitting? (\textbf{\ja{あみもの}}, \emph{amimono}.) To dye? (\textbf{\ja{そめる}}, \emph{someru}.) Embroidery? (\textbf{\ja{ししゅう}}, \emph{shishū}.) To whittle? (\textbf{\ja{けずる}}, \emph{kezuru}.)
\end{tcolorbox}
